\documentclass[11pt]{article}

% Plantilla común para material de estudio y ayudantías.
% Documento A4, una columna, fuente Computer Modern de 11 puntos
% y márgenes de 2.5 cm por lado, siguiendo el formato usado
% en los desarrollos de Física Matemática II.

\usepackage[utf8]{inputenc}
\usepackage[T1]{fontenc}
\usepackage[spanish]{babel}
\usepackage{amsmath,amssymb,mathtools,bm,mathrsfs,graphicx,array,enumitem,caption,float,microtype}
\usepackage[dvipsnames]{xcolor}
\usepackage[a4paper,margin=2.5cm]{geometry}
\usepackage{tikz}
\usetikzlibrary{arrows.meta,calc,patterns,babel,decorations.pathreplacing,decorations.pathmorphing,positioning}
\usepackage[
    colorlinks=true,
    linkcolor=red,
    citecolor=green,
    urlcolor=red
]{hyperref}
\spanishdecimal{.}

\setlength{\parindent}{0pt}
\setlength{\parskip}{0.45em}
\setlength{\abovedisplayskip}{6pt}
\setlength{\belowdisplayskip}{6pt}
\setlength{\abovedisplayshortskip}{4pt}
\setlength{\belowdisplayshortskip}{4pt}
\setlist[itemize]{topsep=4pt,itemsep=2pt,leftmargin=1.8em}
\setlist[enumerate]{topsep=4pt,itemsep=2pt,leftmargin=1.8em}
\captionsetup{font=small,labelfont=bf}
\allowdisplaybreaks

% Colores usados en las figuras.
\definecolor{headinggray}{RGB}{55,55,55}
\definecolor{rulegray}{RGB}{150,150,150}
\definecolor{bluevec}{RGB}{29,78,216}
\definecolor{orangevec}{RGB}{180,83,9}
\definecolor{redvec}{RGB}{185,28,28}
\definecolor{greenvec}{RGB}{21,128,61}
\definecolor{fieldblue}{RGB}{30,90,180}
\definecolor{currentorange}{RGB}{195,95,20}
\definecolor{inducedgreen}{RGB}{20,120,80}
\definecolor{wiregray}{RGB}{55,55,55}
\definecolor{waveblue}{RGB}{29,78,216}
\definecolor{waveorange}{RGB}{180,83,9}
\definecolor{wavegreen}{RGB}{21,128,61}
\definecolor{wavered}{RGB}{185,28,28}

% Encabezado homogéneo.
\newcommand{\encabezado}[3]{%
{\Large\bfseries #1\par}
\vspace{2pt}
{\normalsize #2\hfill #3\par}
\vspace{6pt}\hrule\vspace{8pt}}

\newcommand{\problema}[2]{%
\vspace{0.55em}
\noindent\textbf{#1.}\enspace{\small #2}\par
\vspace{0.3em}}

\newcommand{\inciso}[1]{\vspace{0.35em}\noindent\textbf{(#1)}\enspace}

% Comandos auxiliares.
\newcommand{\dd}{\,\mathrm{d}}
\newcommand{\ii}{\mathrm{i}}
\newcommand{\e}{\mathrm{e}}
\newcommand{\sen}{\operatorname{sen}}
\newcommand{\grad}{\nabla}
\newcommand{\laplaciano}{\nabla^2}
\newcommand{\R}{\mathbb{R}}

\newcommand{\soluciontitulo}{%
  \par\medskip
  \noindent\textbf{Solución.}\par
  \smallskip
}

\newcommand{\resultadofinal}[1]{%
  \par\medskip
  \noindent\textbf{Resultado final.} #1\par
}

\newcommand{\paso}[1]{%
  \par\medskip
  \noindent\textbf{#1}\par
  \smallskip
}

% Entornos sin recuadro: se usan solo para integrar observaciones breves en la redacción.
\newenvironment{comentario}{\par\medskip\noindent}{\par\medskip}
\newenvironment{alerta}{\par\medskip\noindent\textbf{Observación.}\enspace}{\par\medskip}

\newcommand{\Bentra}[2]{%
    \draw[fieldblue,line width=0.55pt] (#1,#2) circle (0.11);
    \draw[fieldblue,line width=0.55pt] ($(#1,#2)+(-0.065,-0.065)$)--($(#1,#2)+(0.065,0.065)$);
    \draw[fieldblue,line width=0.55pt] ($(#1,#2)+(-0.065,0.065)$)--($(#1,#2)+(0.065,-0.065)$);
}
\newcommand{\Bsale}[2]{%
    \draw[fieldblue,line width=0.55pt] (#1,#2) circle (0.11);
    \fill[fieldblue] (#1,#2) circle (0.03);
}

\tikzset{
    >=Latex,
    eje/.style={->,thick,headinggray},
    onda/.style={very thick,waveblue},
    ondados/.style={very thick,waveorange},
    ondaaux/.style={thick,wavegreen},
    vec/.style={->,very thick,wavered},
    vecV/.style={->,very thick,bluevec},
    vecB/.style={->,very thick,orangevec},
    vecF/.style={->,very thick,redvec},
    vecE/.style={->,very thick,greenvec},
    panel/.style={draw=rulegray!80,fill=black!2,rounded corners=2pt},
    cargaPos/.style={circle,draw=bluevec,fill=blue!8,thick,minimum size=8mm,inner sep=0pt},
    cargaNeg/.style={circle,draw=greenvec,fill=green!8,thick,minimum size=8mm,inner sep=0pt},
    guia/.style={densely dashed,rulegray},
    etiqueta/.style={font=\footnotesize},
    nota/.style={draw=rulegray,fill=white,rounded corners=1.5pt,align=left,font=\footnotesize,inner xsep=6pt,inner ysep=5pt,line width=0.35pt}
}

\newcommand{\Fsale}[2]{%
    \draw[redvec,very thick] (#1,#2) circle (0.16);
    \fill[redvec] (#1,#2) circle (0.04);
}
\newcommand{\Fentra}[2]{%
    \draw[redvec,very thick] (#1,#2) circle (0.16);
    \draw[redvec,very thick] ($(#1,#2)+(-0.095,-0.095)$)--($(#1,#2)+(0.095,0.095)$);
    \draw[redvec,very thick] ($(#1,#2)+(-0.095,0.095)$)--($(#1,#2)+(0.095,-0.095)$);
}

\begin{document}

\encabezado{Potencial eléctrico}{Campos y Ondas 510245}{J.R., F.B.}
\shorthandoff{<>}


\section*{Problema 1.}
En una regi\'on del espacio existe un campo el\'ectrico uniforme
\begin{equation*}
    \vec E=(5\hat{x}-3\hat{y}+8\hat{z})\,\text{V/m}.
\end{equation*}
Dados los puntos $A(4,0,0)$, $B(4,8,0)$, $C(0,0,0)$ y $D(10,-4,6)$ en cm, calcular $V_{AB}$ y $V_{DC}$. En este problema usaremos la convenci\'on
\begin{equation*}
    V_{AB}=V_A-V_B,
    \qquad
    V_{DC}=V_D-V_C.
\end{equation*}

\par\medskip\noindent En un campo el\'ectrico uniforme, la diferencia de potencial entre dos puntos se obtiene a partir del producto escalar entre $\vec E$ y el vector desplazamiento entre esos puntos.\par\medskip 

\soluciontitulo

\textbf{Datos del problema.}

El campo el\'ectrico est\'a dado por
\begin{equation*}
    \vec E=(5\hat{x}-3\hat{y}+8\hat{z})\,\text{V/m}.
\end{equation*}

Las posiciones de los puntos, escritas en forma vectorial, son
\begin{align*}
    \vec r_A&=4\,\hat{x}\,\text{cm}+0\,\hat{y}\,\text{cm}+0\,\hat{z}\,\text{cm},\\
    \vec r_B&=4\,\hat{x}\,\text{cm}+8\,\hat{y}\,\text{cm}+0\,\hat{z}\,\text{cm},\\
    \vec r_C&=0\,\hat{x}\,\text{cm}+0\,\hat{y}\,\text{cm}+0\,\hat{z}\,\text{cm},\\
    \vec r_D&=10\,\hat{x}\,\text{cm}-4\,\hat{y}\,\text{cm}+6\,\hat{z}\,\text{cm}.
\end{align*}

Como debemos trabajar en unidades SI, pasamos todo a metros:
\begin{align*}
    \vec r_A&=0.04\,\hat{x}\,\text{m}+0\,\hat{y}\,\text{m}+0\,\hat{z}\,\text{m},\\
    \vec r_B&=0.04\,\hat{x}\,\text{m}+0.08\,\hat{y}\,\text{m}+0\,\hat{z}\,\text{m},\\
    \vec r_C&=0\,\hat{x}\,\text{m}+0\,\hat{y}\,\text{m}+0\,\hat{z}\,\text{m},\\
    \vec r_D&=0.10\,\hat{x}\,\text{m}-0.04\,\hat{y}\,\text{m}+0.06\,\hat{z}\,\text{m}.
\end{align*}

\textbf{Modelo f\'isico.}

La diferencia de potencial entre dos puntos se calcula mediante
\begin{equation*}
    V(P_2)-V(P_1)=-\int_{P_1}^{P_2}\vec E\cdot d\vec \ell.
\end{equation*}
Como el campo el\'ectrico es uniforme, la integral se reduce a
\begin{equation*}
    V(P_2)-V(P_1)=-\vec E\cdot(\vec r_2-\vec r_1).
\end{equation*}

\textbf{C\'alculo de $V_{AB.}$}

Primero construimos el vector desplazamiento desde $A$ hasta $B$:
\begin{align*}
    \Delta \vec r_{AB}
    &=\vec r_B-\vec r_A\\
    &=(0.04\,\hat{x}+0.08\,\hat{y}+0\,\hat{z})\,\text{m}
      -(0.04\,\hat{x}+0\,\hat{y}+0\,\hat{z})\,\text{m}\\
    &=0\,\hat{x}\,\text{m}+0.08\,\hat{y}\,\text{m}+0\,\hat{z}\,\text{m}.
\end{align*}

Entonces,
\begin{equation*}
    V_B-V_A=-\vec E\cdot\Delta \vec r_{AB}.
\end{equation*}

Sustituyendo los datos de forma expl\'icita:
\begin{align*}
    V_B-V_A
    &=-\left(5\hat{x}-3\hat{y}+8\hat{z}\right)\cdot
      \left(0\hat{x}+0.08\hat{y}+0\hat{z}\right)\\
    &=-\left[5(0)+(-3)(0.08)+8(0)\right]\\
    &=-(-0.24)\\
    &=0.24\,\text{V}.
\end{align*}

Como el problema pide
\begin{equation*}
    V_{AB}=V_A-V_B,
\end{equation*}
obtenemos finalmente
\begin{equation*}
    \boxed{V_{AB}=-0.24\,\text{V}.}
\end{equation*}

\textbf{C\'alculo de $V_{DC.}$}

Ahora construimos el vector desplazamiento desde $D$ hasta $C$:
\begin{align*}
    \Delta \vec r_{DC}
    &=\vec r_C-\vec r_D\\
    &=(0\,\hat{x}+0\,\hat{y}+0\,\hat{z})\,\text{m}
      -(0.10\,\hat{x}-0.04\,\hat{y}+0.06\,\hat{z})\,\text{m}\\
    &=-0.10\,\hat{x}\,\text{m}+0.04\,\hat{y}\,\text{m}-0.06\,\hat{z}\,\text{m}.
\end{align*}

Entonces,
\begin{equation*}
    V_C-V_D=-\vec E\cdot\Delta \vec r_{DC}.
\end{equation*}

Sustituyendo:
\begin{align*}
    V_C-V_D
    &=-\left(5\hat{x}-3\hat{y}+8\hat{z}\right)\cdot
      \left(-0.10\hat{x}+0.04\hat{y}-0.06\hat{z}\right)\\
    &=-\left[5(-0.10)+(-3)(0.04)+8(-0.06)\right]\\
    &=-\left[-0.50-0.12-0.48\right]\\
    &=-(-1.10)\\
    &=1.10\,\text{V}.
\end{align*}

Como el problema pide
\begin{equation*}
    V_{DC}=V_D-V_C,
\end{equation*}
resulta
\begin{equation*}
    \boxed{V_{DC}=-1.10\,\text{V}.}
\end{equation*}

\textbf{Resultado final.}
\begin{equation*}
    \boxed{V_{AB}=-0.24\,\text{V}},
    \qquad
    \boxed{V_{DC}=-1.10\,\text{V}}.
\end{equation*}

\newpage

\section*{Problema 2.}
Dos cargas puntuales $q_1=4.0\times 10^{-8}\,\text{C}$ y $q_2=-3.0\times 10^{-8}\,\text{C}$ est\'an separadas $10\,\text{cm}$. De la figura se lee que
\begin{equation*}
    r_{1A}=8\,\text{cm},
    \qquad
    r_{2A}=6\,\text{cm},
    \qquad
    r_{1M}=r_{2M}=5\,\text{cm},
\end{equation*}
donde $M$ es el punto medio entre ambas cargas.

Se pide:
\begin{enumerate}[label=(\alph*)]
    \item El potencial en el punto $A$ y en el punto $M$.
    \item El trabajo hecho por un agente externo para transportar una carga de $25\times 10^{-9}\,\text{C}$ desde $A$ hasta $M$.
\end{enumerate}

\begin{figure}[h!]
    \centering
    \begingroup
    \shorthandoff{<>}
    \begin{tikzpicture}[scale=0.82, line cap=round, line join=round, ]
        % puntos principales
        \coordinate (Q1) at (0,0);
        \coordinate (Q2) at (10,0);
        \coordinate (M)  at (5,0);
        \coordinate (A)  at (6.4,4.8);

        % lados del triángulo
        \draw[thick] (Q1) -- (A) -- (Q2);
        \draw[thick] (Q1) -- (Q2);

        % punto medio
        \fill (M) circle (2.2pt);
        \node[below=4pt] at (M) {$M$};

        % cargas y punto A
        \fill[blue!70!black] (Q1) circle (2.5pt);
        \fill[red!70!black] (Q2) circle (2.5pt);
        \fill (A) circle (2.2pt);

        \node[below=5pt] at (Q1) {$q_1$};
        \node[below=5pt] at (Q2) {$q_2$};
        \node[above=4pt] at (A) {$A$};

        % etiquetas de distancia
        \node[left] at ($(Q1)!0.5!(A)+(-0.25,0.05)$) {$8\,\text{cm}$};
        \node[right] at ($(Q2)!0.5!(A)+(0.25,0.05)$) {$6\,\text{cm}$};
        \node[below] at ($(Q1)!0.5!(Q2)$) {$10\,\text{cm}$};
    \end{tikzpicture}
    \endgroup
    \caption{Geometría usada en el problema 2.}
\end{figure}

\par\medskip\noindent El potencial el\'ectrico es una magnitud escalar. Por eso, para varias cargas puntuales se suman directamente los t\'erminos $kq/r$ con su signo correspondiente.\par\medskip 

\soluciontitulo

\textbf{Datos del problema.}
\begin{align*}
    q_1&=4.0\times 10^{-8}\,\text{C},
    & q_2&=-3.0\times 10^{-8}\,\text{C},\\
    r_{1A}&=0.08\,\text{m},
    & r_{2A}&=0.06\,\text{m},\\
    r_{1M}&=r_{2M}=0.05\,\text{m},
    & q_0&=25\times 10^{-9}\,\text{C}.
\end{align*}
Adem\'as,
\begin{equation*}
    k=9.0\times 10^9\,\text{N m}^2/\text{C}^2.
\end{equation*}

\textbf{Modelo f\'isico.}

El potencial debido a varias cargas puntuales es
\begin{equation*}
    V(P)=k\sum_i \frac{q_i}{r_i}.
\end{equation*}
Para el trabajo del agente externo, bajo un traslado cuasiest\'atico,
\begin{equation*}
    W_{\text{ext}}=\Delta U=q_0(V_f-V_i).
\end{equation*}

\textbf{Ecuaciones de trabajo.}

En $A$:
\begin{equation*}
    V_A=k\left(\frac{q_1}{r_{1A}}+\frac{q_2}{r_{2A}}\right).
\end{equation*}
En $M$:
\begin{equation*}
    V_M=k\left(\frac{q_1}{r_{1M}}+\frac{q_2}{r_{2M}}\right).
\end{equation*}
Finalmente,
\begin{equation*}
    W_{\text{ext}}=q_0(V_M-V_A).
\end{equation*}

\textbf{C\'alculo de $V_A$.}

Sustituyendo,
\begin{align*}
    V_A
    &=9.0\times 10^9\left(\frac{4.0\times 10^{-8}}{0.08}+\frac{-3.0\times 10^{-8}}{0.06}\right).
\end{align*}
Ahora bien,
\begin{equation*}
    \frac{4.0\times 10^{-8}}{0.08}=5.0\times 10^{-7},
    \qquad
    \frac{3.0\times 10^{-8}}{0.06}=5.0\times 10^{-7}.
\end{equation*}
Luego,
\begin{equation*}
    V_A=9.0\times 10^9\big(5.0\times 10^{-7}-5.0\times 10^{-7}\big)=0.
\end{equation*}
Por consiguiente,
\begin{equation*}
    \boxed{V_A=0\,\text{V}}.
\end{equation*}

\textbf{C\'alculo de $V_M$.}

Como $M$ es el punto medio, ambas distancias valen $0.05\,\text{m}$, as\'i que
\begin{align*}
    V_M
    &=9.0\times 10^9\left(\frac{4.0\times 10^{-8}}{0.05}+\frac{-3.0\times 10^{-8}}{0.05}\right)\\
    &=9.0\times 10^9\left(\frac{1.0\times 10^{-8}}{0.05}\right)\\
    &=9.0\times 10^9\times 2.0\times 10^{-7}\\
    &=1.8\times 10^3\,\text{V}.
\end{align*}
Por tanto,
\begin{equation*}
    \boxed{V_M=1800\,\text{V}}.
\end{equation*}

\textbf{C\'alculo del trabajo externo.}

Tomando como estado inicial $A$ y como estado final $M$,
\begin{align*}
    W_{\text{ext}}
    &=q_0(V_M-V_A)\\
    &=(25\times 10^{-9})(1800-0)\\
    &=4.5\times 10^{-5}\,\text{J}.
\end{align*}
En consecuencia,
\begin{equation*}
    \boxed{W_{\text{ext}}=4.5\times 10^{-5}\,\text{J}}.
\end{equation*}

\newpage

\end{document}